% ============================================================================
% THE FLAGSHIP — a machine-checked proof of Chen's theorem
% Target: arXiv (math.NT, cross math.LO/cs.LO) -> Forum of Mathematics, Pi
% Voice: docs/writeup/pre-outline.md §2½ (the ten rules + rule 11).
% Convention: every stated result carries its Lean name via \leanname{};
% the back-matter table maps statement -> declaration -> axiom audit.
% Skeleton 2026-07-15; [TODO] marks drafting sites. No local LaTeX toolchain
% on the build machine — compile checks happen on the author's side/arXiv.
% ============================================================================
\documentclass[11pt]{amsart}
\usepackage[margin=1.1in]{geometry}
\usepackage{amsmath,amssymb,amsthm}
\usepackage{booktabs}
\usepackage[colorlinks=true,linkcolor=blue!60!black,citecolor=blue!60!black,urlcolor=blue!60!black]{hyperref}
\usepackage{microtype}

% --- The Lean-anchor convention -------------------------------------------
% Every formal statement in the paper carries its corpus declaration name.
\newcommand{\leanname}[1]{\par\noindent{\small\texttt{#1}}}
\newcommand{\leaninline}[1]{{\small\texttt{#1}}}
% ---------------------------------------------------------------------------

\newtheorem{theorem}{Theorem}[section]
\newtheorem{corollary}[theorem]{Corollary}
\newtheorem{definition}[theorem]{Definition}
\theoremstyle{remark}
\newtheorem{remark}[theorem]{Remark}

\title{A machine-checked proof of Chen's theorem}
\author{Jason Hickey}
% Personal capacity; standard disclaimer in the acknowledgments (approval
% positioning per the 2026-07-15 decision). Claude's role: Methods, not
% authorship, per journal AI policies; the commit-level attribution is part
% of the artifact.
\date{\today}

\begin{document}

\begin{abstract}
% The §2½ calibration sample, to be refined at drafting. The seven-sentence
% arc: problem -> difficulty -> recourse -> honest state of the art ->
% "here we present" -> mechanism -> quantified result. Rule 3: every strong
% adjective paid by a number in this paragraph.
The correctness of classical analytic number theory rests on proofs whose
critical details are checked by referee trust; verifying them formally has
been notoriously laborious, and for the hardest genre---sieve theory with
explicit constants---essentially never done. Here we present a
machine-checked proof of Chen's theorem: there are infinitely many primes
$p$ such that $p+2$ is prime or a product of two primes, unconditional,
verified by the Lean kernel from three standard axioms. The proof was
constructed by an AI-orchestrated method in days of wall-clock time, on a
corpus of roughly 190{,}000 lines that includes the first formalizations of
an unconditional Siegel--Walfisz theorem, a Bombieri--Vinogradov theorem,
and the switched linear sieve. Formalization produced a complete ledger of
78 corrections to the working presentation, none in the celebrated ideas
and essentially all at the interfaces between subsystems. Every defect was
caught before a wrong proof could be built on it. We state the effective
intermediate results, describe the architecture, and measure where
correctness actually lives in a classical proof. To demonstrate that the
corpus is an instrument rather than a monument, we then posed six
pre-registered exploratory questions against it and report every outcome
with its cost. Two challenges remain:
[TODO: pick the two -- candidates: effective thresholds (the tower), and
the parity frontier].
\end{abstract}

\maketitle

% ============================================================================
\section{Introduction}
% Rule 1: problem-first. No history-of-the-field throat-clearing.
% ============================================================================

[TODO: opening paragraph -- the problem: the details of sieve-theoretic
proofs are trusted, not checked; what "checked" would even mean at this
scale. One paragraph.]

[TODO: "Here we present" paragraph -- the headline theorem, the corpus, the
method in one sentence each. The numbers: 78/0; $\sim$190k lines; the
day-count; axioms.]

\begin{theorem}[Chen's theorem, machine-checked]\label{thm:chen}
There are infinitely many primes $p$ such that $p+2$ is prime or a product
of two primes.
\leanname{Salt.Chen.chen\_headline :
  \{p | p.Prime \wedge IsP2 2 (p + 2)\}.Infinite ---
  axioms: [propext, Classical.choice, Quot.sound]}
\end{theorem}

[TODO: what "machine-checked" means here -- the kernel, the axiom audit,
what the reader must trust (the statement rendering + three axioms +
the toolchain) and what they need not trust (everything else). One
paragraph, plain.]

[TODO: the contribution list, enumerated, each one sentence:
(1) the theorem and its stack (Thm~\ref{thm:chen}, \S\ref{sec:stack});
(2) the effective intermediates (\S\ref{sec:effective});
(3) the ledger and its taxonomy (\S\ref{sec:ledger});
(4) the delimitation theorems (\S\ref{sec:atlas});
(5) the method, summarized here and detailed in the companion
(\S\ref{sec:method}); the artifact (\S\ref{sec:artifact}).]

[TODO: the honesty paragraph, per §1½: no new headline theorem is claimed;
what IS new (the formalizations-first, the minor new statements, the
measurement); the Kepler precedent \cite{HalesKepler} for the genre.]

% ============================================================================
\section{The theorem and its stack}\label{sec:stack}
% ============================================================================

[TODO: the mathematical statement unpacked: IsP2, the history one paragraph
(Chen 1966/1973 \cite{Chen1973}, the switching idea), what unconditional
means here (no GRH, no Siegel-zero case split surviving to the statement).]

\subsection{The analytic spine}
[TODO: one subsection per layer, each with the formal statement + Lean
anchor + a mechanism paragraph:]

\begin{theorem}[Unconditional Siegel--Walfisz]\label{thm:sw}
[TODO: the statement with explicit constants.]
\leanname{Salt.SW.siegelWalfisz\_holds}
\end{theorem}

\begin{theorem}[Bombieri--Vinogradov]\label{thm:bv}
[TODO: the dispersion form; the maxDiscrepancy convention.]
\leanname{Salt.BV.psi\_BV\_of\_siegelWalfisz'}
\end{theorem}

\begin{theorem}[The windowed bilinear Bombieri--Vinogradov]\label{thm:wbv}
[TODO: from Salt/Chen/WindowedBVStatement.lean -- the literature-facing
form; the Q4 verdict: the conjunction (bilinear + sharp window + explicit
constants + the block-prime $\sqrt{D}$ level) is unrecorded. State both the
general form and the clean corollary.]
\leanname{Salt.Chen.windowed\_bilinear\_BV\_sqrtD}
\end{theorem}

[TODO: the sieve layer: the linear sieve at the Rosser--Iwaniec chain, the
switch, the weighted count at $\bar c/2$, the razor with
$M = 0.012151$; each anchored.]

\subsection{The operating point}
[TODO: the witness: $z = x^{1/8}$, $y = x^{1/3}$, the $W$-trick modulus,
the tower threshold -- and the honest sentence about effectivity: the
threshold is astronomically large; the theorem is about existence, the
effective-threshold problem is Challenge 1 in \S\ref{sec:challenges}.]

% ============================================================================
\section{Effective intermediate results}\label{sec:effective}
% The dividend section: statements a classical reader can cite without Lean.
% ============================================================================

[TODO: the table of effective statements with constants: the SW constants,
the BV level and constants, the windowed-BV Kerr form, the razor margin,
the $\bar c$ certificate (600 panels, error $3.18\times 10^{-8}$), the
mass-ledger values. Each row: statement, constant, Lean anchor.]

% ============================================================================
\section{The ledger: where correctness lives}\label{sec:ledger}
% The anatomy section, condensed. The Notices companion carries the long
% form; here: the data + the three case studies.
% ============================================================================

[TODO: the setup paragraph: statements were frozen before proof; every
failed freeze is a recorded catch; 78 catches, 0 proofs on wrong
statements; the discovery-mechanism split (executor 47\%, gate 31\%,
inventory 12\%, kernel 0) -- the kernel confirms, never discovers.]

[TODO: the taxonomy table (from docs/writeup/taxonomy.md): categories x
phase; the headline correlation: integration-phase catches are 91\%
interface-shaped; build-phase 60\% local (80\% excluding embedded
mini-integrations).]

\subsection{Case study: the one-sentence lemma}
[TODO: "by dyadic decomposition we may assume..." -> the pricing layer:
catches \#59--\#78 in three paragraphs.]

\subsection{Case study: shape-correct but uninhabitable}
[TODO: catch \#65 -- the H-package tear; typechecking vs joint
satisfiability; the W-trick repair.]

\subsection{Case study: the modal failure}
[TODO: carrier mismatch -- the most common integration defect (11 of 23);
the hcount seam as the exhibit.]

% ============================================================================
\section{The instrument turned around: six pre-registered questions}
\label{sec:atlas}
% The demo section (JYH framing 2026-07-15: "this is a platform for
% exploratory math; to illustrate we posed 6 problems"). The frame is ONE
% paragraph + ONE table; the section's mass stays the theorems themselves.
% The full process anatomy (gates, tears, token economics) belongs to the
% method paper; the exploration ledger in the repo is the raw artifact
% both papers cite.
% ============================================================================

[TODO: the demo frame, one paragraph: to test whether the corpus is a
research instrument rather than a monument, we posed six exploratory
questions in a pre-registered commit BEFORE exploring (the hash,
disclosed here as the tamper-evident timestamp) and report every outcome
with its cost -- including questions that returned obstruction maps
rather than theorems. The pre-registration is what converts "some nice
corollaries" into a controlled demonstration: no cherry-picking, negative
results reported.]

[TODO: the table: question / outcome class (theorem; statement +
literature verdict; obstruction map; reuse measurement) / cost. Six rows;
costs from docs/writeup/token-data.md, outcomes from the exploration
ledger at sprint close.]

The theorems the questions returned:

\begin{theorem}[The $k=2$ Maynard bound]\label{thm:m2}
For every continuous $F$ on the simplex,
$J_1(F)+J_2(F) \le 2\log 2 \cdot I(F)$; in particular no admissible weight
crosses the twin gate.
\leanname{Salt.TwinBar.twin\_bar, no\_twin\_weight}
\end{theorem}

\begin{theorem}[Enlargement does not help at $k=2$]\label{thm:m2e}
[TODO: the $(1+\delta)$ statement with $\delta_0 = 1/\log 2 - 1$; the
constrained form and Tao's open problem \cite{Polymath8bVII}; the numeric
map of $c(\delta)$.]
\leanname{Salt.TwinBar.no\_twin\_weight\_enlarged,
no\_twin\_weight\_constrained}
\end{theorem}

\begin{theorem}[$M_3 < 2$]\label{thm:m3}
[TODO: the $(3/2)\log 3$ statement, unconditional.]
\leanname{Salt.TwinBar.three\_bar, no\_triple\_weight}
\end{theorem}

[TODO if landed by drafting: the parity wall (the Selberg-witness interface
theorem) and the interface extraction gaps\_le\_twelve\_of\_hasLevel with
the $\theta$-caveat stated plainly.]

% ============================================================================
\section{The method, in brief}\label{sec:method}
% One page. The companion paper carries the full treatment.
% ============================================================================

[TODO: the roles (designer/executors/kernel), the freeze-gate-wave-ceremony
loop, witness-first refinement (one paragraph on premise latency), the
stats that matter here: 468 commits / 8.44 days; 178/178 first-attempt
landings with the honest caveat; the endgame economics (13.0M output
tokens, $\approx$650k/catch). Cite the companion.]

% ============================================================================
\section{The artifact}\label{sec:artifact}
% ============================================================================

[TODO: the repo, the tagged release, clone-and-continue (the build in two
commands), the audit instructions (\#print axioms), the machine-readable
ledger, the open-problems board. "The proof is not attached to this paper;
this paper is attached to the proof."]

% ============================================================================
\section{Limitations, and two challenges}\label{sec:challenges}
% Rule 7: enumerated, actionable, at the end.
% ============================================================================

[TODO: the limitations paragraph: the tower threshold (ineffective-in-
practice operating point); the constants are explicit but not competitive;
the statement-rendering trust boundary; what the 78-catch measurement can
and cannot generalize to.]

\textbf{Challenge 1.} [TODO: effective thresholds -- bring the operating
point down from the tower; what blocks it (the $w_0$-window guards).]

\textbf{Challenge 2.} [TODO: the parity frontier -- TwinB\_min as the
precisely-stated door; what a formal assault would need.]

% ============================================================================
\section*{Acknowledgments}
% ============================================================================
[TODO: the personal-capacity disclaimer (the approved formula); the Claude
role statement (orchestration + proof construction per the Methods;
commit-level attribution in the artifact; not an author per journal
policy); mathlib and its community; the classical sources (BJS, Tao's
expositions) whose renderings this work audits.]

% ============================================================================
% Back matter: the statement -> declaration -> audit table
% ============================================================================
\appendix
\section{The formal-statement index}\label{app:index}

[TODO: the full table. Columns: paper statement, Lean declaration, file,
axiom audit result. Generated from Salt/Chen/All.lean's audit block +
Salt/TwinBar/All.lean; keep in sync by script at release.]

\begin{table}[h]
\centering\small
\begin{tabular}{lll}
\toprule
Statement & Declaration & Axioms \\
\midrule
Thm~\ref{thm:chen} & \leaninline{Salt.Chen.chen\_headline} &
  \leaninline{propext, Classical.choice, Quot.sound} \\
Thm~\ref{thm:sw} & \leaninline{Salt.SW.siegelWalfisz\_holds} & (same) \\
Thm~\ref{thm:bv} & \leaninline{Salt.BV.psi\_BV\_of\_siegelWalfisz'} & (same) \\
Thm~\ref{thm:wbv} & \leaninline{Salt.Chen.windowed\_bilinear\_BV\_sqrtD} & (same) \\
Thm~\ref{thm:m2} & \leaninline{Salt.TwinBar.no\_twin\_weight} & (same) \\
Thm~\ref{thm:m3} & \leaninline{Salt.TwinBar.no\_triple\_weight} & (same) \\
\bottomrule
\end{tabular}
\caption{[TODO: the complete index at release.]}
\end{table}

\begin{thebibliography}{99}
\bibitem{Chen1973} J.~R. Chen, \emph{On the representation of a larger even
integer as the sum of a prime and the product of at most two primes},
Sci.\ Sinica \textbf{16} (1973), 157--176.
\bibitem{BJS} [TODO: Bordignon--Johnston--Starichkova, the explicit Chen
source, arXiv:2207.09452.]
\bibitem{TaoExp} [TODO: Tao's expository treatment(s) used for the
transcription.]
\bibitem{HalesKepler} T.~Hales et al., \emph{A formal proof of the Kepler
conjecture}, Forum of Mathematics, Pi \textbf{5} (2017), e2.
\bibitem{Gonthier} G.~Gonthier, \emph{Formal proof---the four-color
theorem}, Notices Amer.\ Math.\ Soc.\ \textbf{55} (2008), 1382--1393.
\bibitem{Polymath8b} D.~H.~J. Polymath, \emph{Variants of the Selberg
sieve, and bounded intervals containing many primes}, Res.\ Math.\ Sci.\
\textbf{1} (2014).
\bibitem{Polymath8bVII} [TODO: the Polymath8b VII post/paper for the k=2
GEH problem (Q5a/Tao's suspicion).]
\bibitem{AkbaryHambrook} [TODO: Akbary--Hambrook, the explicit BV,
arXiv:1309.2730.]
\bibitem{mathlib} The mathlib community, \emph{The Lean mathematical
library}, CPP 2020.
\bibitem{companion} [TODO: the method companion paper, in preparation.]
\end{thebibliography}

\end{document}
